%% ----------------------------------------------------------------------------
%%  §4.6  Discontinuità e prolungamento per continuità
%% ----------------------------------------------------------------------------
\section{Discontinuità}
\label{sec:discontinuita}

La funzione
\begin{equation}\label{eq:segno}
  f(x)=\frac{\abs x}{x}=
  \begin{cases}
    1 & \text{se } x>0,\\
    -1 & \text{se } x<0
  \end{cases}
\end{equation}
è continua per $x\neq0$ (in un intorno di ogni $x_{0}\neq0$ è costante, e il limite è
una proprietà locale, \cref{lem:localita}), ma non è continua se $x=0$. Il grafico di
questa funzione presenta per $x=0$ un salto, appunto una \emph{discontinuità}
(\cref{fig:modulo-su-x}). In particolare $f$ non è continua nel punto $x_{0}=0$ perché
non è definita in tal punto, cioè perché non esiste il valore $f(x_{0})=f(0)$.

Estendiamo allora $f$ anche a $x_{0}=0$ con un valore $\ell\in\R$, considerando cioè la
nuova funzione
\begin{equation}\label{eq:segno-esteso}
  \tilde f(x)=
  \begin{cases}
    \dfrac{\abs x}{x} & \text{se } x\neq0,\\[1ex]
    \ell & \text{se } x=0 .
  \end{cases}
\end{equation}
La funzione $\tilde f$ è definita anche nel punto $x_{0}=0$, ma non è continua in tale
punto, perché non esiste il limite per $x\to0$ di $\tilde f(x)$: precisamente, il limite
destro è diverso dal limite sinistro \eqref{eq:modulo-su-x}. Tutto ciò accade
\emph{qualunque sia} il valore $\ell$ scelto in \eqref{eq:segno-esteso}: $\tilde f$ è
un'estensione non continua (o \emph{prolungamento non continuo}) della funzione $f$, e
la discontinuità di $f$ nel punto $x_{0}=0$ si dice \emph{non eliminabile}.

Viceversa, la funzione studiata nel \cref{sec:premessa-funzioni},
$f(x)=\frac{\sen x}{x}$, non è continua in $x_{0}=0$ (perché non è definita), ma è
possibile \emph{prolungarla per continuità} mediante la funzione
\begin{equation}\label{eq:sinc}
  \tilde f(x)=
  \begin{cases}
    \dfrac{\sen x}{x} & \text{se } x\neq0,\\[1ex]
    1 & \text{se } x=0 .
  \end{cases}
\end{equation}
A causa del limite notevole \eqref{eq:lim-senx-x}, $\tilde f$ è continua anche nel punto
$x_{0}=0$. Il grafico di $\tilde f$ si ottiene «completando» il disegno della
\cref{fig:senx-x} con l'ulteriore punto di coordinate $(0,1)$, che riempie il tondino
vuoto.

\subsection{Classificazione delle discontinuità}

Sia $f$ una funzione definita in $A$ e $x_{0}$ un punto di $A$, di accumulazione per $A$.
Se $f$ non è continua in $x_{0}$ si dice che $x_{0}$ è un punto di \emph{discontinuità}
per $f$. Le discontinuità si classificano nel modo seguente
(\cref{fig:discontinuita}).

\begin{definizione}[Tipi di discontinuità]\label{def:discontinuita}\index{discontinuità}
\begin{enumerate}[label=\textup{(\alph*)}]
  \item La funzione $f$ presenta in $x_{0}$ una \emph{discontinuità eliminabile} se
  esiste finito il limite di $f(x)$ per $x\to x_{0}$ e risulta
  \begin{equation}\label{eq:disc-eliminabile}
    \lim_{x\to x_{0}}f(x)\neq f(x_{0}) .
  \end{equation}
  \item La funzione $f$ presenta in $x_{0}$ una \emph{discontinuità di prima specie} (o
  \emph{di salto}) se esistono finiti i limiti destro e sinistro di $f(x)$ in $x_{0}$ e
  si ha
  \begin{equation}\label{eq:disc-prima-specie}
    \lim_{x\to x_{0}^{-}}f(x)\neq\lim_{x\to x_{0}^{+}}f(x);
  \end{equation}
  la differenza tra il limite destro e il limite sinistro si chiama \emph{salto} di $f$
  in $x_{0}$.
  \item La funzione $f$ presenta in $x_{0}$ una \emph{discontinuità di seconda specie}
  se almeno uno dei due limiti
  \begin{equation}\label{eq:disc-seconda-specie}
    \lim_{x\to x_{0}^{-}}f(x), \qquad \lim_{x\to x_{0}^{+}}f(x)
  \end{equation}
  non esiste oppure è infinito.
\end{enumerate}
\end{definizione}

La classificazione è completa: se $f$ è discontinua in $x_{0}$ (interno al dominio) e
non siamo nel caso (c), i due limiti unilaterali esistono finiti; se sono diversi siamo
nel caso (b), se sono uguali il limite esiste (\cref{prop:dx-sx}) ed è diverso da
$f(x_{0})$, e siamo nel caso (a). Per abuso di linguaggio, le stesse parole si usano
anche per un punto $x_{0}$ di accumulazione che \emph{non} appartiene al dominio, come
$x_{0}=0$ per $\frac{\sen x}{x}$ (singolarità eliminabile) o per $\frac{\abs x}{x}$
(salto).

Nel caso (a), posto $\ell=\lim_{x\to x_{0}}f(x)$, la funzione
\begin{equation}\label{eq:elimina}
  \tilde f(x)=
  \begin{cases}
    f(x) & \text{se } x\in A\setminus\{x_{0}\},\\
    \ell & \text{se } x=x_{0}
  \end{cases}
\end{equation}
risulta continua nel punto $x_{0}$: la discontinuità si «elimina» correggendo un solo
valore. Nei casi (b) e (c) nessuna modifica del valore in $x_{0}$ rende $f$ continua.

\begin{figure}[htb]
\centering
\includegraphics[width=\textwidth]{discontinuita}
\caption{I tre tipi di discontinuità. (a) La funzione vale $\frac{\sen x}{x}$ per
$x\neq0$ e $\frac12$ in $0$: il limite esiste, ma è diverso dal valore. (b) I limiti
destro e sinistro esistono finiti ma sono diversi. (c) La funzione $e^{1/x}$ ha limite
sinistro $0$ e limite destro $+\infty$ (asse $y$ in scala logaritmica per $y>1$)
(\texttt{python/fig\_discontinuita.py}).}
\label{fig:discontinuita}
\end{figure}

\begin{esempio}\label{es:parte-intera-continuita}
La funzione $f(x)=\lfloor x\rfloor$, \emph{parte intera} di $x$
(\cref{prop:parte-intera}, rappresentata nella \cref{fig:parte-intera}), presenta
discontinuità di prima specie in corrispondenza a ogni $k\in\Z$, mentre è continua per
ogni $x\in\R\setminus\Z$. Infatti, se $x_{0}\notin\Z$, esiste $k\in\Z$ con
$k<x_{0}<k+1$, e nell'intorno $(k,k+1)$ di $x_{0}$ la funzione è costante, uguale a $k$.
Se invece $x_{0}=k\in\Z$, si ha $\lfloor x\rfloor=k-1$ per $k-1\le x<k$ e
$\lfloor x\rfloor=k$ per $k\le x<k+1$; quindi
\[
  \lim_{x\to k^{-}}\lfloor x\rfloor=k-1, \qquad \lim_{x\to k^{+}}\lfloor x\rfloor=k=\lfloor k\rfloor:
\]
il salto vale $1$, e la funzione è continua da destra ma non da sinistra.
\end{esempio}

\begin{esempio}\label{es:seconda-specie}
\begin{enumerate}[label=\textup{(\roman*)}]
  \item La funzione $f(x)=\sen\frac1x$ per $x\neq0$, prolungata con un qualsiasi valore
  in $0$, ha in $0$ una discontinuità di seconda specie: il limite destro non esiste
  (lo si vede con le successioni dell'\cref{es:sen-1-x}, che sono positive).
  \item La funzione $f(x)=\frac1x$ per $x\neq0$, $f(0)=0$, ha in $0$ una discontinuità di
  seconda specie: i limiti unilaterali sono infiniti (\cref{es:modulo-su-x}).
  \item La funzione $e^{1/x}$ ha in $0$ una singolarità di seconda specie: per $x\to0^{+}$
  si ha $\frac1x\to+\infty$ ed $e^{y}\to+\infty$ per $y\to+\infty$, quindi, per il
  \cref{teo:limite-composta} (con $y_{0}=+\infty$, caso in cui l'ipotesi $g(x)\neq y_{0}$
  è automatica), $e^{1/x}\to+\infty$; per $x\to0^{-}$, invece, $\frac1x\to-\infty$ ed
  $e^{1/x}\to0$ (\cref{fig:discontinuita}(c)).
\end{enumerate}
\end{esempio}

\begin{osservazione}[Una funzione discontinua ovunque]\label{oss:dirichlet}\index{funzione!di Dirichlet}
La \emph{funzione di Dirichlet}
\[
  D(x)=
  \begin{cases}
    1 & \text{se } x\in\Q,\\
    0 & \text{se } x\in\R\setminus\Q
  \end{cases}
\]
è discontinua in ogni punto, con discontinuità di seconda specie. Fissato $x_{0}\in\R$,
per la densità dei razionali e degli irrazionali (\cref{teo:densita-Q},
\cref{cor:densita-irrazionali}) in ogni intervallo $(x_{0},x_{0}+\frac1n)$ esistono un
razionale $q_{n}$ e un irrazionale $r_{n}$; le successioni $q_{n}$ ed $r_{n}$ tendono a
$x_{0}$ da destra, ma $D(q_{n})=1\to1$ e $D(r_{n})=0\to0$, quindi il limite destro non
esiste (\cref{cor:non-esistenza}). Peter Gustav Lejeune Dirichlet la propose nel 1829,
nel suo lavoro sulle serie di Fourier, come esempio di funzione per cui l'integrale, e
con esso la sua analisi, perde significato; nel 1837 formulò esplicitamente la nozione di
funzione come corrispondenza arbitraria, priva di qualunque regolarità. Questi esempi
contribuirono a imporre la definizione moderna di funzione.
\end{osservazione}

\subsection{Il prolungamento per continuità}

\begin{definizione}[Prolungamento per continuità]\label{def:prolungamento}\index{prolungamento per continuità}
Sia $A$ un intervallo (o un'unione finita di intervalli), $x_{0}\in A$, e sia $f$ una
funzione definita in $A\setminus\{x_{0}\}$. Se esiste finito il limite
\begin{equation}\label{eq:prolungamento-limite}
  \lim_{x\to x_{0}}f(x)=\ell,
\end{equation}
la funzione $\tilde f$ definita in $A$ da
\begin{equation}\label{eq:prolungamento}
  \tilde f(x)=
  \begin{cases}
    f(x) & \text{se } x\in A\setminus\{x_{0}\},\\
    \ell & \text{se } x=x_{0}
  \end{cases}
\end{equation}
è detta \emph{prolungamento per continuità di $f$ in $x_{0}$}; per costruzione $\tilde f$
risulta continua in $x_{0}$. Se poi $f$ è continua in $A\setminus\{x_{0}\}$, allora
$\tilde f$, continua su tutto l'insieme $A$, è detta \emph{prolungamento per continuità
di $f$ su $A$}.
\end{definizione}

Il prolungamento per continuità, quando esiste, è unico, perché il limite è unico.
La \eqref{eq:sinc} è il prolungamento per continuità di $\frac{\sen x}{x}$ su $\R$; la
funzione che si ottiene, detta \emph{seno cardinale} e indicata con $\operatorname{sinc}x$,
è di uso quotidiano nell'elaborazione dei segnali: è, a meno di costanti, la trasformata
di Fourier di un impulso rettangolare, e descrive ad esempio il diagramma di
diffrazione di un'apertura o di un'antenna. Nei programmi il prolungamento va scritto
esplicitamente.

\begin{python}[Il prolungamento per continuità nei programmi]
import math
import numpy as np

try:
    print(math.sin(0.0) / 0.0)
except ZeroDivisionError as err:
    print("sen(0)/0 ->", type(err).__name__, ":", err)

def sinc(x):
    """Prolungamento per continuità di sen(x)/x: vale 1 in x = 0."""
    return 1.0 if x == 0 else math.sin(x) / x

print(sinc(0.0), sinc(1e-9), sinc(math.pi / 2))
# NumPy usa la versione normalizzata sen(pi x)/(pi x), estesa in 0
print(np.sinc(0.0), np.sinc(0.5), np.sinc(1.0))
\end{python}
\begin{risultato}
sen(0)/0 -> ZeroDivisionError : float division by zero
1.0 1.0 0.6366197723675814
1.0 0.6366197723675814 3.8981718325193755e-17
\end{risultato}

L'ultimo valore dovrebbe essere $\sen\pi/\pi=0$: è il solito residuo di arrotondamento
dovuto alla rappresentazione approssimata di $\pi$.

\begin{idea}[L'onda d'urto: una discontinuità di prima specie]
In un flusso supersonico può formarsi un'onda d'urto normale: attraversandola, la
pressione, la densità e la temperatura aumentano bruscamente e la velocità diminuisce.
Nel modello di fluido non viscoso (equazioni di Eulero) l'urto è una vera discontinuità
di prima specie: lungo l'asse della corrente la pressione vale $p_{1}$ a monte e $p_{2}$
a valle, con
\[
  \frac{p_{2}}{p_{1}}=1+\frac{2\gamma}{\gamma+1}\bigl(M_{1}^{2}-1\bigr),
\]
dove $M_{1}$ è il numero di Mach a monte e $\gamma=1{,}4$ per l'aria; con $M_{1}=2$ il
salto porta la pressione a $4{,}5$ volte il valore di monte. Nella realtà la viscosità
«smussa» l'urto in uno strato continuo spesso pochi cammini liberi medi delle molecole,
frazioni di micrometro al livello del mare: alla scala di un'ala, una discontinuità a
tutti gli effetti. Nei codici CFD con tecniche di \emph{shock capturing} lo stesso salto
viene invece distribuito su poche celle di calcolo: la soluzione numerica è continua e
approssima, raffinando la griglia, una soluzione discontinua.
\end{idea}
